\documentclass[11pt]{article}

\usepackage[margin=1in]{geometry}
\usepackage{amsmath,amssymb,amsthm,mathtools}
\usepackage{enumitem}
\usepackage{microtype}
\usepackage{xcolor}
\usepackage{hyperref}
\usepackage{url}
\usepackage{booktabs}
\usepackage{array}

\hypersetup{
  colorlinks=true,
  linkcolor=blue!55!black,
  citecolor=blue!55!black,
  urlcolor=blue!55!black
}

\newtheorem{theorem}{Theorem}[section]
\newtheorem{proposition}[theorem]{Proposition}
\newtheorem{lemma}[theorem]{Lemma}
\newtheorem{corollary}[theorem]{Corollary}
\newtheorem{remark}[theorem]{Remark}
\newtheorem{definition}[theorem]{Definition}
\newtheorem{convention}[theorem]{Convention}

\newcommand{\Hom}{\operatorname{Hom}}
\newcommand{\homn}{\operatorname{hom}}
\newcommand{\tr}{\operatorname{Tr}}
\newcommand{\EE}{\mathbb{E}}
\newcommand{\RR}{\mathbb{R}}
\newcommand{\NN}{\mathbb{N}}
\newcommand{\1}{\mathbf{1}}
\newcommand{\e}{\mathrm{e}}
\newcommand{\Ct}{\mathcal{C}}
\newcommand{\wtOmega}{\widetilde{\Omega}}
\newcommand{\dd}{\,\mathrm{d}}

\title{Even-Anchored Interpolation for Odd Cycles\\
       and Consequences for Domination Exponents}
\author{Detailed working note}
\date{August 2026}

\begin{document}
\maketitle

\begin{abstract}
Let $T$ be a finite simple target graph.  We prove that whenever $b$ is even,
$n,m$ are odd, and $2\le b<n<m$, one has
\[
  t(\Ct_b,T)^{m-n}\,t(C_m,T)^{n-b}
  \ge t(C_n,T)^{m-b},
\]
where $\Ct_2:=K_2$ and $\Ct_b:=C_b$ for $b\ge4$.
The proof is obtained by rooting all three closed walks at a vertex, applying
Sa\u{g}lam's same-parity power inequality to a nonnegative vector of the form
$A^{b/2}e_x$, and then summing by H\"older's inequality.  In particular, this
proves the cycle inequality posed as Problem~6 in Blekherman--Raymond--Razborov--Wei.
Combining the edge-anchored case $b=2$ with their pruning argument yields
\[
 C(C_{2k+1},C_{2\ell+1})
 \le
 \frac{2k(2k-2\ell+1)-1}
      {2\ell(2k-2\ell+1)-1}
 \qquad (k>\ell),
\]
thereby extending their second upper-bound formula to the full parameter range.
Every reduction after Sa\u{g}lam's theorem is written out explicitly, including
zero cases, normalization from homomorphism numbers to densities, the pruning
constants, the elimination of the edge density, and the tensor-power limit.
A final section records a Lean-oriented dependency graph and theorem statements.
\end{abstract}

\tableofcontents

\section{Scope, notation, and logical dependencies}

All graphs in this note are finite, simple, undirected, and loopless.  A target
graph is denoted by $T$, its vertex set by $V(T)$, and its number of vertices by
\[
  N:=|V(T)|.
\]
We always assume $N\ge1$.

For finite graphs $F,T$, let
\[
  \homn(F,T):=|\Hom(F,T)|
\]
be the number of graph homomorphisms $F\to T$, and let
\[
  t(F,T):=\frac{\homn(F,T)}{N^{|V(F)|}}
\]
be the corresponding homomorphism density.

\begin{convention}[The symbol $\Ct_2$]
For the purpose of a uniform trace formula, define
\[
  \Ct_2:=K_2,
  \qquad
  \Ct_r:=C_r\quad(r\ge3).
\]
Thus $\Ct_2$ is not a two-edge multigraph.  This convention is valid for
finite simple targets because the entries of their adjacency matrices are
$0$ or $1$.
\end{convention}

The proof has one non-elementary external input.

\begin{theorem}[Sa\u{g}lam's same-parity power inequality]\label{thm:saglam}
Let $\Omega$ be a finite set.  Let $S\in\RR^{\Omega\times\Omega}$ be symmetric
and entrywise nonnegative, and let $u,v\in\RR^\Omega$ be entrywise nonnegative
unit vectors.  If $K,R$ are positive integers satisfying
\[
  K\ge R,
  \qquad
  K\equiv R\pmod 2,
\]
then
\[
  \langle v,S^K u\rangle^R
  \ge
  \langle v,S^R u\rangle^K.
\]
\end{theorem}

This is Theorem~1.3 of Sa\u{g}lam~\cite{Saglam2018}.  We use only the special
case $u=v$.  No regularity, stochasticity, or positive-semidefiniteness is
assumed: symmetry and entrywise nonnegativity are the relevant hypotheses.

\medskip
\noindent
\textbf{Dependency boundary.}
The exact cycle-interpolation theorem below is completely proved from
Theorem~\ref{thm:saglam}, finite-dimensional matrix algebra, and H\"older's
inequality.  The domination-exponent corollary additionally uses a pruning
argument, which is reproduced in Section~\ref{sec:pruning}, and the elementary
even-cycle estimate proved in Lemma~\ref{lem:even-cycle-edge}.  Thus a formal
proof assistant development may isolate Theorem~\ref{thm:saglam} as the only
substantial imported theorem.  Every use of it below records the matrix,
vectors, exponents, parity checks, zero-denominator case, and denominator
clearing explicitly.

\begin{remark}[Novelty status]
The inequality proved here is exactly the inequality asked as Problem~6 in
version~2 of Blekherman--Raymond--Razborov--Wei~\cite{BRRW2025}.  This note
records a proof and its consequences; it does not make a priority claim.
Independent verification should precede any claim of novelty.
\end{remark}

\section{Adjacency matrices and rooted closed walks}

Let $A$ be the adjacency matrix of $T$, indexed by $V(T)$.  Thus
\[
  A_{xy}=\begin{cases}
  1,&\{x,y\}\in E(T),\\
  0,&\text{otherwise}.
  \end{cases}
\]
The matrix $A$ is symmetric and entrywise nonnegative.  For $x\in V(T)$, let
$e_x$ be the standard basis vector at $x$.  For every integer $r\ge0$, define
\[
  c_r(x):=(A^r)_{xx}=\langle e_x,A^r e_x\rangle.
\]
For $r\ge1$, $c_r(x)$ is the number of closed walks of length $r$ in $T$ that
start and end at $x$.

\begin{lemma}[Trace and homomorphism counts]\label{lem:trace-hom}
For every integer $r\ge3$,
\[
  \homn(C_r,T)=\tr(A^r)=\sum_{x\in V(T)}c_r(x).
\]
Moreover,
\[
  \homn(K_2,T)=\tr(A^2)=\sum_{x\in V(T)}c_2(x)=2|E(T)|.
\]
Consequently, for every $r\ge2$,
\[
  \homn(\Ct_r,T)=\tr(A^r).
\]
\end{lemma}

\begin{proof}
Fix $r\ge3$, and write the vertices of $C_r$ cyclically as
$0,1,\ldots,r-1$.  Expanding the $(x,x)$ entry of $A^r$ gives
\[
  (A^r)_{xx}
  =
  \sum_{x_1,\ldots,x_{r-1}\in V(T)}
  A_{x x_1}A_{x_1x_2}\cdots A_{x_{r-2}x_{r-1}}A_{x_{r-1}x}.
\]
Because every factor is either $0$ or $1$, a summand equals $1$ exactly when
\[
  0\mapsto x,
  \quad
  1\mapsto x_1,
  \quad\ldots,\quad
  r-1\mapsto x_{r-1}
\]
defines a homomorphism $C_r\to T$.  Therefore $c_r(x)$ counts precisely the
homomorphisms for which vertex $0$ is mapped to $x$.  Summing over $x$ counts
every homomorphism exactly once, proving the first identity.

For $r=2$,
\[
  (A^2)_{xx}=\sum_{y\in V(T)}A_{xy}A_{yx}
  =\sum_y A_{xy}^2
  =\sum_y A_{xy}
  =\deg_T(x),
\]
where symmetry and $A_{xy}\in\{0,1\}$ were used.  Hence
\[
  \tr(A^2)=\sum_x\deg_T(x)=2|E(T)|.
\]
A homomorphism $K_2\to T$ is an ordered adjacent pair, so its number is also
$2|E(T)|$.
\end{proof}

The next identity is the reason an even cycle can be used as an anchor.

\begin{lemma}[Moving an even anchor to the endpoints]\label{lem:anchor-vector}
Let $b\ge2$ be even, let $x\in V(T)$, and set
\[
  w_x:=A^{b/2}e_x.
\]
Then $w_x$ is entrywise nonnegative and, for every integer $r\ge b$,
\[
  \|w_x\|_2^2=c_b(x),
  \qquad
  \langle w_x,A^{r-b}w_x\rangle=c_r(x).
\]
\end{lemma}

\begin{proof}
Every entry of every nonnegative integer power of the entrywise nonnegative
matrix $A$ is nonnegative, so $w_x\ge0$ coordinatewise.  Since $A$ is
symmetric, so is $A^{b/2}$, and therefore
\[
  \|w_x\|_2^2
  =\langle A^{b/2}e_x,A^{b/2}e_x\rangle
  =\langle e_x,A^{b/2}A^{b/2}e_x\rangle
  =\langle e_x,A^b e_x\rangle
  =c_b(x).
\]
Similarly, for $r\ge b$,
\[
\begin{aligned}
  \langle w_x,A^{r-b}w_x\rangle
  &=\langle A^{b/2}e_x,A^{r-b}A^{b/2}e_x\rangle\\
  &=\langle e_x,A^{b/2}A^{r-b}A^{b/2}e_x\rangle\\
  &=\langle e_x,A^r e_x\rangle
  =c_r(x).
\end{aligned}
\]
All powers commute because they are powers of the same matrix.
\end{proof}

\section{The rooted interpolation inequality}

\begin{lemma}[Rooted even-anchor interpolation]\label{lem:rooted-interpolation}
Let $b,n,m$ be integers such that
\[
  2\le b<n<m,
  \qquad
  b\text{ is even},
  \qquad
  n,m\text{ are odd}.
\]
Then, for every $x\in V(T)$,
\[
  c_b(x)^{m-n}\,c_m(x)^{n-b}
  \ge
  c_n(x)^{m-b}.
  \tag{3.1}\label{eq:rooted-interpolation}
\]
\end{lemma}

\begin{proof}
Fix $x\in V(T)$ and let $w_x=A^{b/2}e_x$ as in
Lemma~\ref{lem:anchor-vector}.  We split into two cases.

\smallskip
\noindent
\emph{Case 1: $c_b(x)=0$.}
By Lemma~\ref{lem:anchor-vector},
\[
  \|w_x\|_2^2=c_b(x)=0.
\]
A vector in a real inner-product space has norm zero only when it is the zero
vector, so $w_x=0$.  Again by Lemma~\ref{lem:anchor-vector},
\[
  c_n(x)=\langle w_x,A^{n-b}w_x\rangle=0,
  \qquad
  c_m(x)=\langle w_x,A^{m-b}w_x\rangle=0.
\]
All three exponents in \eqref{eq:rooted-interpolation} are positive, because
$b<n<m$.  Thus both sides are zero, and the inequality holds.

\smallskip
\noindent
\emph{Case 2: $c_b(x)>0$.}
Define
\[
  u_x:=\frac{w_x}{\|w_x\|_2}.
\]
Then $u_x$ is an entrywise nonnegative unit vector.  Put
\[
  K:=m-b,
  \qquad
  R:=n-b.
\]
Since $b$ is even and $m,n$ are odd, both $K$ and $R$ are positive odd
integers.  Moreover $K>R$.  The hypotheses of
Theorem~\ref{thm:saglam} therefore hold with
\[
  S=A,
  \qquad
  u=v=u_x.
\]
Hence
\[
  \langle u_x,A^{m-b}u_x\rangle^{n-b}
  \ge
  \langle u_x,A^{n-b}u_x\rangle^{m-b}.
  \tag{3.2}\label{eq:saglam-applied}
\]
By Lemma~\ref{lem:anchor-vector} and
$\|w_x\|_2^2=c_b(x)$,
\[
  \langle u_x,A^{r-b}u_x\rangle
  =\frac{\langle w_x,A^{r-b}w_x\rangle}{\|w_x\|_2^2}
  =\frac{c_r(x)}{c_b(x)}
  \qquad(r\in\{n,m\}).
\]
Substitution into \eqref{eq:saglam-applied} gives
\[
  \left(\frac{c_m(x)}{c_b(x)}\right)^{n-b}
  \ge
  \left(\frac{c_n(x)}{c_b(x)}\right)^{m-b}.
\]
Multiplying both sides by the positive number $c_b(x)^{m-b}$ yields
\[
  c_b(x)^{(m-b)-(n-b)}c_m(x)^{n-b}
  \ge c_n(x)^{m-b}.
\]
Since $(m-b)-(n-b)=m-n$, this is exactly
\eqref{eq:rooted-interpolation}.
\end{proof}

\section{Why the regular case is easy}
\label{sec:regular}

The ordinary regularity assumption corresponds to the special anchor $b=2$.
Indeed,
\[
  c_2(x)=\deg_T(x).
\]
For $b>2$, ordinary regularity does not generally make $c_b(x)$ independent
of $x$; that stronger property is usually called walk-regularity at length
$b$.  Thus the regular-to-general interpretation below should first be read
for the edge-anchored inequality.

\begin{proposition}[Edge-anchored interpolation for a regular target]
\label{prop:regular-edge-anchor}
Let $2<n<m$ be odd, and let $T$ be a $d$-regular graph on $N$ vertices.  Then
\[
  t(K_2,T)^{m-n}\,t(C_m,T)^{n-2}
  \ge
  t(C_n,T)^{m-2}.
  \tag{4.1}\label{eq:regular-density}
\]
\end{proposition}

\begin{proof}
Set
\[
  \alpha:=\frac{m-n}{m-2},
  \qquad
  \beta:=\frac{n-2}{m-2}.
\]
Because $2<n<m$, one has
\[
  0<\alpha<1,
  \qquad
  0<\beta<1,
  \qquad
  \alpha+\beta=1.
\]
Lemma~\ref{lem:rooted-interpolation}, with $b=2$, gives
\[
  d^{m-n}c_m(x)^{n-2}\ge c_n(x)^{m-2}.
\]
All quantities are nonnegative, so taking the $(m-2)$-th root gives
\[
  c_n(x)\le d^\alpha c_m(x)^\beta.
\]
Summing over $x$,
\[
  \homn(C_n,T)
  =\sum_xc_n(x)
  \le d^\alpha\sum_x c_m(x)^\beta.
\]
Since $0<\beta<1$, H\"older's inequality, equivalently the power-mean
inequality, gives
\[
  \sum_xc_m(x)^\beta
  \le
  N^{1-\beta}\left(\sum_xc_m(x)\right)^\beta
  =N^\alpha\homn(C_m,T)^\beta.
\]
Therefore
\[
  \homn(C_n,T)
  \le
  (Nd)^\alpha\homn(C_m,T)^\beta.
\]
For a $d$-regular graph,
\[
  \homn(K_2,T)=2|E(T)|=Nd.
\]
Thus
\[
  \homn(C_n,T)
  \le
  \homn(K_2,T)^\alpha\homn(C_m,T)^\beta.
\]
Raising to the power $m-2$ gives
\[
  \homn(K_2,T)^{m-n}\homn(C_m,T)^{n-2}
  \ge
  \homn(C_n,T)^{m-2}.
  \tag{4.2}\label{eq:regular-hom}
\]
Finally, the total number of source vertices on the two sides agrees:
\[
  2(m-n)+m(n-2)=n(m-2).
\]
Dividing \eqref{eq:regular-hom} by $N^{n(m-2)}$ therefore gives
\eqref{eq:regular-density}.
\end{proof}

The proof shows exactly what ordinary regularity contributed: it made the
rooted anchor $c_2(x)=d$ constant, so it could be pulled out of the sum.  In
the general case, the same rooted inequality remains true, and a two-factor
H\"older inequality absorbs the variation of $d(x)$ exactly.  There is no
uncontrolled additive error and therefore no need for a Schur complement in
this problem.

\section{Removing regularity by H\"older's inequality}

We first isolate the precise finite H\"older statement used in the summation.

\begin{lemma}[Integer-exponent form of two-factor H\"older]
\label{lem:integer-holder}
Let $I$ be a finite set.  Let $P,Q$ be positive integers and put $R=P+Q$.
Suppose $a_i,b_i,c_i\ge0$ for every $i\in I$ and
\[
  a_i^R\le b_i^P c_i^Q
  \qquad(i\in I).
  \tag{5.1}\label{eq:pointwise-holder-hyp}
\]
Then
\[
  \left(\sum_{i\in I}a_i\right)^R
  \le
  \left(\sum_{i\in I}b_i\right)^P
  \left(\sum_{i\in I}c_i\right)^Q.
  \tag{5.2}\label{eq:integer-holder}
\]
\end{lemma}

\begin{proof}
Set
\[
  \alpha:=\frac{P}{R},
  \qquad
  \beta:=\frac{Q}{R}.
\]
Then $\alpha,\beta\in(0,1)$ and $\alpha+\beta=1$.  Since the map
$x\mapsto x^{1/R}$ is increasing on $[0,\infty)$,
\eqref{eq:pointwise-holder-hyp} implies
\[
  a_i\le b_i^\alpha c_i^\beta.
\]
Let
\[
  p_0:=\frac{1}{\alpha}=\frac{R}{P},
  \qquad
  q_0:=\frac{1}{\beta}=\frac{R}{Q}.
\]
Then $p_0,q_0>1$ and $1/p_0+1/q_0=1$.  H\"older's inequality gives
\[
\begin{aligned}
  \sum_i a_i
  &\le \sum_i b_i^\alpha c_i^\beta\\
  &\le
  \left(\sum_i(b_i^\alpha)^{p_0}\right)^{1/p_0}
  \left(\sum_i(c_i^\beta)^{q_0}\right)^{1/q_0}\\
  &=
  \left(\sum_i b_i\right)^\alpha
  \left(\sum_i c_i\right)^\beta.
\end{aligned}
\]
Both sides are nonnegative.  Raising to the positive integer power $R$ gives
\eqref{eq:integer-holder}.
\end{proof}

\begin{theorem}[Even-anchored interpolation for odd cycles]
\label{thm:cycle-interpolation}
Let $b,n,m$ satisfy
\[
  2\le b<n<m,
  \qquad
  b\text{ even},
  \qquad
  n,m\text{ odd}.
\]
Then every finite simple target graph $T$ satisfies
\[
  \homn(\Ct_b,T)^{m-n}\,
  \homn(C_m,T)^{n-b}
  \ge
  \homn(C_n,T)^{m-b}.
  \tag{5.3}\label{eq:main-hom}
\]
Equivalently,
\[
  t(\Ct_b,T)^{m-n}\,
  t(C_m,T)^{n-b}
  \ge
  t(C_n,T)^{m-b}.
  \tag{5.4}\label{eq:main-density}
\]
\end{theorem}

\begin{proof}
For each $x\in V(T)$, Lemma~\ref{lem:rooted-interpolation} gives
\[
  c_n(x)^{m-b}
  \le
  c_b(x)^{m-n}c_m(x)^{n-b}.
\]
Apply Lemma~\ref{lem:integer-holder} with
\[
  I=V(T),
  \qquad
  a_x=c_n(x),
  \qquad
  b_x=c_b(x),
  \qquad
  c_x=c_m(x),
\]
and with
\[
  P=m-n,
  \qquad
  Q=n-b,
  \qquad
  R=P+Q=m-b.
\]
The hypotheses $b<n<m$ ensure that $P,Q$ are positive.  We obtain
\[
  \left(\sum_xc_n(x)\right)^{m-b}
  \le
  \left(\sum_xc_b(x)\right)^{m-n}
  \left(\sum_xc_m(x)\right)^{n-b}.
\]
By Lemma~\ref{lem:trace-hom}, this is exactly
\eqref{eq:main-hom}.

It remains to verify the equivalence with the density formulation.  Since
$|V(\Ct_b)|=b$, $|V(C_n)|=n$, and $|V(C_m)|=m$, substituting
\[
  \homn(F,T)=N^{|V(F)|}t(F,T)
\]
into \eqref{eq:main-hom} introduces the power
\[
  N^{b(m-n)+m(n-b)}
\]
on the left and the power
\[
  N^{n(m-b)}
\]
on the right.  These exponents agree because
\[
\begin{aligned}
  b(m-n)+m(n-b)
  &=bm-bn+mn-mb\\
  &=mn-bn\\
  &=n(m-b).
\end{aligned}
\]
Since $N>0$, the common power cancels, yielding
\eqref{eq:main-density}.
\end{proof}

\begin{remark}[The exact regular-to-general mechanism]
For $b=2$, the pointwise inequality is
\[
  \deg_T(x)^{m-n}c_m(x)^{n-2}
  \ge c_n(x)^{m-2}.
\]
In a regular graph, $\deg_T(x)$ is constant and can be pulled out of the sum.
In an arbitrary graph, H\"older's inequality replaces that step by
\[
  \sum_x\deg_T(x)^\alpha c_m(x)^\beta
  \le
  \left(\sum_x\deg_T(x)\right)^\alpha
  \left(\sum_xc_m(x)\right)^\beta.
\]
Thus degree irregularity is not bounded perturbatively; it is integrated
exactly.  This is why a Schur-complement correction is unnecessary here.
\end{remark}

\section{The cycle inequality posed as Problem 6}

\begin{corollary}[Affirmative answer to Problem 6 of \cite{BRRW2025}]
\label{cor:problem6}
Let $i>j\ge1$.  Then every finite simple graph $T$ satisfies
\[
  t(C_{2j},T)^2
  t(C_{2i+1},T)^{2i-1-2j}
  \ge
  t(C_{2i-1},T)^{2i+1-2j}.
  \tag{6.1}\label{eq:problem6}
\]
Here, as in \cite{BRRW2025}, $C_2$ means $K_2$.
\end{corollary}

\begin{proof}
Apply Theorem~\ref{thm:cycle-interpolation} with
\[
  b=2j,
  \qquad
  n=2i-1,
  \qquad
  m=2i+1.
\]
The condition $i>j$ implies
\[
  2j\le2i-2<2i-1,
\]
so $b<n<m$.  Also, $b$ is even while $n,m$ are odd.  The three exponents in
\eqref{eq:main-density} become
\[
  m-n=2,
\]
\[
  n-b=(2i-1)-2j=2i-1-2j,
\]
and
\[
  m-b=(2i+1)-2j=2i+1-2j.
\]
Substitution gives exactly \eqref{eq:problem6}.
\end{proof}

\begin{remark}[Tropicalization consequence]
Blekherman--Raymond--Razborov--Wei state that an affirmative solution of their
Problem~6 is equivalent to their proposed description of the tropicalization
of the homomorphism-number profile of all cycles; see their Lemma~6.2.
Corollary~\ref{cor:problem6}, combined with that equivalence, supplies the
missing implication.  We do not reproduce their separate proof of the
equivalence here.
\end{remark}

\section{Domination exponents}

\begin{definition}[Homomorphism-density domination exponent]
For graphs $G,H$, define
\[
  C(G,H)
  :=
  \min\bigl\{c\in\RR:
  t(G,T)\ge t(H,T)^c\text{ for every finite simple }T\bigr\},
\]
whenever the set is nonempty.
\end{definition}

We next derive the domination-exponent bound announced in the abstract.  The
argument uses the edge-anchored case of
Theorem~\ref{thm:cycle-interpolation}, an explicit pruning reduction, and an
elementary even-cycle estimate.

\subsection{An even-cycle lower bound}

\begin{lemma}[Even cycles dominate the appropriate edge power]
\label{lem:even-cycle-edge}
Let $r\ge2$.  Then every finite simple graph $T$ satisfies
\[
  t(C_{2r},T)\ge t(K_2,T)^{2r}.
  \tag{7.1}\label{eq:even-cycle-edge}
\]
\end{lemma}

\begin{proof}
Let $A$ be the adjacency matrix of $T$, and let $p=t(K_2,T)$.  Let
$\lambda_{\max}$ be the largest eigenvalue of $A$.  Since $A$ is real
symmetric, the Rayleigh--Ritz principle and the vector $\1\in\RR^{V(T)}$ give
\[
  \lambda_{\max}
  \ge
  \frac{\langle\1,A\1\rangle}{\langle\1,\1\rangle}.
\]
Now
\[
  \langle\1,A\1\rangle=\homn(K_2,T)=N^2p,
  \qquad
  \langle\1,\1\rangle=N,
\]
so
\[
  \lambda_{\max}\ge Np.
\]
Let $\lambda_1,\ldots,\lambda_N$ be all eigenvalues of $A$, repeated with
multiplicity.  Because $2r$ is even,
\[
  \homn(C_{2r},T)
  =\tr(A^{2r})
  =\sum_{a=1}^N\lambda_a^{2r}
  \ge\lambda_{\max}^{2r}
  \ge (Np)^{2r}.
\]
Dividing by $N^{2r}$ proves \eqref{eq:even-cycle-edge}.
\end{proof}

\subsection{A precise pruning reduction}
\label{sec:pruning}

For an integer $q\ge3$, give $C_q$ the cyclic vertex set
$\mathbb{Z}/q\mathbb{Z}$ and distinguish the oriented edge $(0,1)$.  For an
oriented target edge $(u,v)$, define
\[
  h_q^T(u,v)
  :=
  \bigl|\{\phi\in\Hom(C_q,T):\phi(0)=u,\ \phi(1)=v\}\bigr|.
  \tag{7.2}\label{eq:rooted-edge-count}
\]
Let
\[
  \vec E(T):=\{(u,v)\in V(T)^2:\{u,v\}\in E(T)\}
\]
be the set of oriented edges.

\begin{lemma}[Basic identities for edge-rooted cycle counts]
\label{lem:edge-rooted-identities}
For every $q\ge3$,
\[
  \sum_{(u,v)\in\vec E(T)}h_q^T(u,v)=\homn(C_q,T).
  \tag{7.3}\label{eq:edge-rooted-sum}
\]
Moreover,
\[
  h_q^T(u,v)=h_q^T(v,u)
  \qquad((u,v)\in\vec E(T)).
  \tag{7.4}\label{eq:orientation-symmetry}
\]
\end{lemma}

\begin{proof}
Every homomorphism $\phi:C_q\to T$ maps the distinguished oriented edge
$(0,1)$ to exactly one oriented edge $(u,v)$ of $T$.  Partitioning
$\Hom(C_q,T)$ according to this image proves
\eqref{eq:edge-rooted-sum}.

The map $a\mapsto1-a$ on $\mathbb{Z}/q\mathbb{Z}$ is an automorphism of $C_q$
that interchanges $0$ and $1$.  Precomposition with this automorphism is a
bijection between the homomorphisms counted by $h_q^T(u,v)$ and those counted
by $h_q^T(v,u)$.  This proves \eqref{eq:orientation-symmetry}.
\end{proof}

\begin{lemma}[Pruning to uniformly cycle-rich edges]
\label{lem:pruning}
Fix an integer $q\ge3$.  Let $T_0$ be a graph on $N\ge3$ vertices.  Put
\[
  H_0:=\homn(C_q,T_0).
\]
If $H_0>0$, then $T_0$ has a spanning subgraph $T\subseteq T_0$ such that,
writing
\[
  M:=|E(T)|,
  \qquad
  H:=\homn(C_q,T),
\]
one has
\[
  H\ge \e^{-3/2}H_0,
  \tag{7.5}\label{eq:pruning-preservation}
\]
and, for every oriented edge $(u,v)\in\vec E(T)$,
\[
  h_q^T(u,v)
  >
  \frac{H}{8qM\log N}.
  \tag{7.6}\label{eq:uniform-rooted-edge}
\]
Here and below $\log$ is the natural logarithm.
\end{lemma}

\begin{proof}
We construct a descending sequence of spanning subgraphs.  Suppose the current
graph $T_i$ has exactly $i$ edges, and write
\[
  H_i:=\homn(C_q,T_i).
\]
For an edge $e\in E(T_i)$, let $L_i(e)$ be the number of homomorphisms
$\phi:C_q\to T_i$ whose image uses $e$, meaning that at least one source edge
of $C_q$ is mapped to $e$.

If there is an edge $e\in E(T_i)$ satisfying
\[
  L_i(e)\le\frac{H_i}{4i\log N},
  \tag{7.7}\label{eq:pruning-delete-rule}
\]
delete it and call the resulting graph $T_{i-1}$.  A homomorphism is lost
exactly when its image uses $e$, so
\[
  H_{i-1}=H_i-L_i(e)
  \ge H_i\left(1-\frac{1}{4i\log N}\right).
  \tag{7.8}\label{eq:one-pruning-step}
\]
If no edge satisfies \eqref{eq:pruning-delete-rule}, stop and call the terminal
graph $T$.  The terminal graph has at least one edge: throughout the process
\eqref{eq:one-pruning-step} keeps $H_i>0$, while an edgeless graph admits no
homomorphism from $C_q$.

Let $M_0:=|E(T_0)|$ and $M:=|E(T)|$.  Iterating
\eqref{eq:one-pruning-step} gives
\[
  H
  \ge
  H_0\prod_{i=M+1}^{M_0}
  \left(1-\frac{1}{4i\log N}\right).
  \tag{7.9}\label{eq:pruning-product}
\]
For $N\ge3$,
\[
  0\le\frac{1}{4i\log N}<\frac12.
\]
The elementary inequality
\[
  \log(1-x)\ge-2x
  \qquad(0\le x\le1/2)
\]
therefore yields
\[
\begin{aligned}
  \log\prod_{i=M+1}^{M_0}
  \left(1-\frac{1}{4i\log N}\right)
  &\ge
  -\frac{1}{2\log N}\sum_{i=1}^{M_0}\frac1i\\
  &\ge
  -\frac{1+\log M_0}{2\log N}.
\end{aligned}
\]
Since $M_0\le N^2/2<N^2$,
\[
  1+\log M_0\le1+2\log N\le3\log N,
\]
where the last inequality uses $\log N\ge1$ for $N\ge3$.  Hence the logarithm
of the product is at least $-3/2$, which proves
\eqref{eq:pruning-preservation}.

It remains to prove \eqref{eq:uniform-rooted-edge}.  At termination, every
edge $e=\{u,v\}$ satisfies
\[
  L_T(e)>\frac{H}{4M\log N}.
  \tag{7.10}\label{eq:terminal-contained-count}
\]
For each $a\in\mathbb{Z}/q\mathbb{Z}$, let $I_a(e)$ be the number of
homomorphisms $\phi:C_q\to T$ for which the source edge $\{a,a+1\}$ is mapped
to $e$.  Every homomorphism counted by $L_T(e)$ contributes to at least one
$I_a(e)$, so
\[
  \sum_{a\in\mathbb{Z}/q\mathbb{Z}}I_a(e)\ge L_T(e).
  \tag{7.11}\label{eq:incidences-dominate-copies}
\]
By rotation symmetry of $C_q$, all $I_a(e)$ are equal.  For $a=0$,
\[
  I_0(e)=h_q^T(u,v)+h_q^T(v,u)=2h_q^T(u,v),
\]
where the second equality is \eqref{eq:orientation-symmetry}.  Thus
\[
  2q\,h_q^T(u,v)
  =\sum_aI_a(e)
  \ge L_T(e)
  >\frac{H}{4M\log N}.
\]
Division by $2q$ proves \eqref{eq:uniform-rooted-edge}.
\end{proof}

\begin{lemma}[Gluing two edge-rooted cycles]
\label{lem:cycle-gluing}
Let $q,h\ge3$, and put
\[
  m:=q+h-2.
\]
Suppose a graph $T$ on $N\ge3$ vertices satisfies
\eqref{eq:uniform-rooted-edge} for $C_q$.  Then, with
\[
  a:=t(C_q,T),
  \qquad
  z:=t(C_h,T),
  \qquad
  p:=t(K_2,T),
  \qquad
  y:=t(C_m,T),
\]
one has
\[
  y>\frac{1}{4q\log N}\frac{az}{p}.
  \tag{7.12}\label{eq:gluing-density}
\]
\end{lemma}

\begin{proof}
The hypothesis \eqref{eq:uniform-rooted-edge} implies $E(T)\ne\varnothing$, so
$p>0$.  We define the source graph explicitly.  Give $C_m$ the cyclic vertex set
$\{0,1,\ldots,m-1\}$ and add the chord $\{0,q-1\}$.  Call the resulting graph
$C_m^+$.  The chord together with
\[
  0,1,\ldots,q-1
\]
forms a $q$-cycle.  The same chord together with the outer path
\[
  q-1,q,\ldots,m-1,0
\]
forms a cycle of length
\[
  1+(m-q)+1=m-q+2=h,
\]
where the first and last $1$ count the chord and the edge $\{m-1,0\}$.
Thus $C_m^+$ is exactly the union of a copy of $C_q$ and a copy of $C_h$
with one distinguished oriented edge and both of its endpoints identified.

Every homomorphism $C_m^+\to T$ is, after forgetting the chord constraint, a
homomorphism $C_m\to T$.  Therefore
\[
  \homn(C_m,T)\ge\homn(C_m^+,T).
  \tag{7.13}\label{eq:forget-chord}
\]
Condition on the oriented image $(u,v)$ of the identified edge.  Once $(u,v)$
is fixed, the extensions over the two cycles are independent, and hence
\[
  \homn(C_m^+,T)
  =
  \sum_{(u,v)\in\vec E(T)}h_q^T(u,v)h_h^T(u,v).
  \tag{7.14}\label{eq:gluing-exact}
\]
Using the uniform lower bound \eqref{eq:uniform-rooted-edge}, followed by
\eqref{eq:edge-rooted-sum}, gives
\[
\begin{aligned}
  \homn(C_m^+,T)
  &>
  \frac{\homn(C_q,T)}{8q|E(T)|\log N}
  \sum_{(u,v)\in\vec E(T)}h_h^T(u,v)\\
  &=
  \frac{\homn(C_q,T)\homn(C_h,T)}
       {8q|E(T)|\log N}.
\end{aligned}
\]
Combining with \eqref{eq:forget-chord}, and using
\[
  |E(T)|=\frac12N^2p,
  \quad
  \homn(C_q,T)=N^q a,
  \quad
  \homn(C_h,T)=N^h z,
  \quad
  \homn(C_m,T)=N^m y,
\]
we obtain
\[
  N^my
  >
  \frac{N^{q+h}az}{4qN^2p\log N}.
\]
Since $q+h=m+2$, the powers of $N$ cancel and
\eqref{eq:gluing-density} follows.
\end{proof}

\subsection{Eliminating the edge density}

Fix integers $k>\ell\ge1$, and abbreviate
\[
  \Delta:=k-\ell,
  \qquad
  n:=2\ell+1,
  \qquad
  m:=2k+1,
  \qquad
  h:=2\Delta+2.
\]
Then
\[
  n+h-2=m.
\]
For a target graph $T$, write
\[
  a:=t(C_n,T),
  \qquad
  y:=t(C_m,T),
  \qquad
  p:=t(K_2,T).
\]

\begin{lemma}[The exact edge-anchored constraint]
\label{lem:edge-anchored-constraint}
For every target graph $T$,
\[
  p^{2\Delta}y^{2\ell-1}
  \ge
  a^{2k-1}.
  \tag{7.15}\label{eq:edge-anchored-constraint}
\]
\end{lemma}

\begin{proof}
Apply Theorem~\ref{thm:cycle-interpolation} with
\[
  b=2,
  \qquad
  n=2\ell+1,
  \qquad
  m=2k+1.
\]
Then
\[
  m-n=2(k-\ell)=2\Delta,
  \qquad
  n-2=2\ell-1,
  \qquad
  m-2=2k-1,
\]
which gives \eqref{eq:edge-anchored-constraint}.
\end{proof}

\begin{lemma}[Polylogarithmic lower bound after pruning]
\label{lem:polylog-combined}
There are constants $c_0>0$ and $\gamma\ge0$, depending only on $k,\ell$, such
that every graph $T_0$ on $N\ge3$ vertices satisfies
\[
  t(C_m,T_0)
  \ge
  c_0(\log N)^{-\gamma}
  t(C_n,T_0)^{\eta},
  \tag{7.16}\label{eq:polylog-final}
\]
where
\[
  \eta
  :=
  \frac{2k(2\Delta+1)-1}
       {2\ell(2\Delta+1)-1}.
  \tag{7.17}\label{eq:eta-definition}
\]
\end{lemma}

\begin{proof}
Put
\[
  a_0:=t(C_n,T_0),
  \qquad
  y_0:=t(C_m,T_0).
\]
If $a_0=0$, then \eqref{eq:polylog-final} is immediate because $y_0\ge0$.
Assume $a_0>0$.  Then $\homn(C_n,T_0)>0$, so
Lemma~\ref{lem:pruning} gives a spanning subgraph $T\subseteq T_0$ satisfying
\eqref{eq:pruning-preservation} and
\eqref{eq:uniform-rooted-edge}.  Write
\[
  a:=t(C_n,T),
  \qquad
  y:=t(C_m,T),
  \qquad
  p:=t(K_2,T).
\]
Because $T$ is spanning, the density normalization uses the same $N$, and
\eqref{eq:pruning-preservation} gives
\[
  a\ge \e^{-3/2}a_0.
  \tag{7.18}\label{eq:a-preserved}
\]
Since $T\subseteq T_0$, every homomorphism to $T$ is also a homomorphism to
$T_0$, and hence
\[
  y_0\ge y.
  \tag{7.19}\label{eq:y-monotone}
\]

Apply Lemma~\ref{lem:cycle-gluing} with $q=n$ and the even integer
$h=2\Delta+2$.  Then apply Lemma~\ref{lem:even-cycle-edge} to $C_h$.  We get
\[
\begin{aligned}
  y
  &>\frac{1}{4n\log N}
      \frac{a\,t(C_h,T)}{p}\\
  &\ge\frac{1}{4n\log N}a p^{h-1}\\
  &=\frac{1}{4n\log N}a p^{2\Delta+1}.
\end{aligned}
\tag{7.20}\label{eq:pruned-lower}
\]
On the other hand, Lemma~\ref{lem:edge-anchored-constraint} gives
\[
  p^{2\Delta}y^{2\ell-1}\ge a^{2k-1}.
  \tag{7.21}\label{eq:exact-constraint-repeated}
\]

To eliminate $p$, introduce the positive integers
\[
  A_0:=2\ell-1,
  \qquad
  B_0:=2\Delta,
  \qquad
  C_0:=2k-1=A_0+B_0,
  \qquad
  D_0:=2\Delta+1=B_0+1.
\]
Equations \eqref{eq:pruned-lower} and
\eqref{eq:exact-constraint-repeated} become
\[
  y>\frac{1}{4n\log N}a p^{D_0},
  \tag{7.22}\label{eq:two-ineq-one}
\]
and
\[
  p^{B_0}y^{A_0}\ge a^{C_0}.
  \tag{7.23}\label{eq:two-ineq-two}
\]
Since $a>0$, \eqref{eq:two-ineq-two} implies $p>0$ and $y>0$.  Raise
\eqref{eq:two-ineq-one} to the power $B_0$:
\[
  y^{B_0}
  >
  (4n\log N)^{-B_0}a^{B_0}p^{B_0D_0}.
  \tag{7.24}\label{eq:raised-pruning}
\]
Raise \eqref{eq:two-ineq-two} to the power $D_0$ and rearrange:
\[
  p^{B_0D_0}
  \ge
  \frac{a^{C_0D_0}}{y^{A_0D_0}}.
  \tag{7.25}\label{eq:p-elimination}
\]
Substitution of \eqref{eq:p-elimination} into
\eqref{eq:raised-pruning} yields
\[
  y^{B_0+A_0D_0}
  >
  (4n\log N)^{-B_0}a^{B_0+C_0D_0}.
  \tag{7.26}\label{eq:combined-power}
\]
Therefore
\[
  y
  >
  (4n)^{-\theta}(\log N)^{-\theta}
  a^{\eta},
  \tag{7.27}\label{eq:combined-root}
\]
where
\[
  \theta:=\frac{B_0}{B_0+A_0D_0}
\]
and
\[
  \eta:=\frac{B_0+C_0D_0}{B_0+A_0D_0}.
\]
We simplify the denominator:
\[
\begin{aligned}
  B_0+A_0D_0
  &=B_0+A_0(B_0+1)\\
  &=(A_0+1)(B_0+1)-1\\
  &=2\ell(2\Delta+1)-1.
\end{aligned}
\]
Similarly,
\[
\begin{aligned}
  B_0+C_0D_0
  &=B_0+(A_0+B_0)(B_0+1)\\
  &=(A_0+B_0+1)(B_0+1)-1\\
  &=2k(2\Delta+1)-1.
\end{aligned}
\]
This proves \eqref{eq:eta-definition}.  Finally combine
\eqref{eq:a-preserved}, \eqref{eq:y-monotone}, and
\eqref{eq:combined-root}:
\[
\begin{aligned}
  y_0
  &\ge y\\
  &>(4n)^{-\theta}(\log N)^{-\theta}a^\eta\\
  &\ge
  (4n)^{-\theta}\e^{-3\eta/2}
  (\log N)^{-\theta}a_0^\eta.
\end{aligned}
\]
Thus \eqref{eq:polylog-final} holds with
\[
  c_0=(4n)^{-\theta}\e^{-3\eta/2},
  \qquad
  \gamma=\theta.
\]
\end{proof}

\subsection{Removing the polylogarithmic loss}

\begin{lemma}[Tensor-power removal of a polylogarithmic factor]
\label{lem:tensor-trick}
Let $H_1,H_2$ be fixed graphs and let $c\ge0$.  Suppose there are constants
$K>0$ and $D\ge0$ such that, for every graph $T$ on $N\ge3$ vertices,
\[
  t(H_1,T)\ge K(\log N)^{-D}t(H_2,T)^c.
  \tag{7.28}\label{eq:tensor-hypothesis}
\]
Then
\[
  t(H_1,T)\ge t(H_2,T)^c
\]
for every graph $T$ having at least three vertices.
\end{lemma}

\begin{proof}
Fix a target graph $T$ with $N\ge3$ vertices.  For $r\ge1$, let
$T^{\times r}$ be its $r$-fold categorical tensor product.  Its vertex set has
size $N^r$.  For every fixed source graph $H$, a map
$V(H)\to V(T)^r$ is a homomorphism to $T^{\times r}$ exactly when each of its
$r$ coordinate maps is a homomorphism to $T$.  Hence
\[
  t(H,T^{\times r})=t(H,T)^r.
  \tag{7.29}\label{eq:tensor-density}
\]
Apply \eqref{eq:tensor-hypothesis} to $T^{\times r}$:
\[
  t(H_1,T)^r
  \ge
  K\bigl(\log(N^r)\bigr)^{-D}t(H_2,T)^{cr}
  =K(r\log N)^{-D}t(H_2,T)^{cr}.
\]
Taking the $r$-th root gives
\[
  t(H_1,T)
  \ge
  K^{1/r}(r\log N)^{-D/r}t(H_2,T)^c.
\]
As $r\to\infty$,
\[
  K^{1/r}\longrightarrow1,
  \qquad
  (r\log N)^{-D/r}
  =\exp\left(-\frac{D\log(r\log N)}{r}\right)
  \longrightarrow1.
\]
Therefore $t(H_1,T)\ge t(H_2,T)^c$.

This proves the assertion for every target having at least three vertices.
\end{proof}

\begin{theorem}[Uniform upper bound for odd-cycle domination exponents]
\label{thm:domination-bound}
For all integers $k>\ell\ge1$,
\[
  C(C_{2k+1},C_{2\ell+1})
  \le
  \frac{2k(2k-2\ell+1)-1}
       {2\ell(2k-2\ell+1)-1}.
  \tag{7.30}\label{eq:domination-bound}
\]
\end{theorem}

\begin{proof}
Lemma~\ref{lem:polylog-combined} gives, for every target $T$ with at least
three vertices,
\[
  t(C_{2k+1},T)
  \ge
  c_0(\log |V(T)|)^{-\gamma}
  t(C_{2\ell+1},T)^\eta
\]
with
\[
  \eta
  =
  \frac{2k(2(k-\ell)+1)-1}
       {2\ell(2(k-\ell)+1)-1}.
\]
Lemma~\ref{lem:tensor-trick} removes the polylogarithmic factor and gives
\[
  t(C_{2k+1},T)\ge t(C_{2\ell+1},T)^\eta
\]
for every such target.

It remains only to consider targets with at most two vertices.  Every finite
simple graph on at most two vertices is bipartite.  A homomorphism from an odd
cycle to a bipartite graph cannot exist: composing such a homomorphism with a
proper two-colouring of the target would give a proper two-colouring of the
odd cycle.  Hence both odd-cycle densities in the preceding inequality are
zero, and the inequality is $0\ge0$.

Finally, $2(k-\ell)+1=2k-2\ell+1$, so $\eta$ is exactly the right-hand side of
\eqref{eq:domination-bound}.  The defining property of the domination exponent
therefore gives the claimed upper bound.
\end{proof}

\begin{corollary}[Adjacent odd cycles]
For every $i\ge2$,
\[
  C(C_{2i+1},C_{2i-1})
  \le
  \frac{6i-1}{6i-7}.
\]
\end{corollary}

\begin{proof}
Set $k=i$ and $\ell=i-1$ in
Theorem~\ref{thm:domination-bound}.  Then
$2k-2\ell+1=3$, so the numerator is $6i-1$ and the denominator is
$6(i-1)-1=6i-7$.
\end{proof}

\subsection{Comparison with the previously stated piecewise upper bound}

Blekherman--Raymond--Razborov--Wei~\cite{BRRW2025} prove the upper bound
\[
 C(C_{2k+1},C_{2\ell+1})
 \le
 \begin{cases}
 \dfrac{2(k+1)}{2\ell+1},&k\le2\ell-1,\\[2mm]
 \dfrac{2k(2k-2\ell+1)-1}
       {2\ell(2k-2\ell+1)-1},&k\ge2\ell.
 \end{cases}
\]
Theorem~\ref{thm:domination-bound} extends the second expression to all
$k>\ell$.  In the former first-branch range, put
\[
  D:=2k-2\ell+1.
\]
A direct common-denominator calculation gives
\[
\begin{aligned}
  \frac{2(k+1)}{2\ell+1}
  -\frac{2kD-1}{2\ell D-1}
  &=\frac{(4\ell-2k-1)D}
  {(2\ell+1)(2\ell D-1)}.
\end{aligned}
\]
If $k\le2\ell-1$, then $4\ell-2k-1\ge1$ and $D>0$, so the improvement is
strict.

Their projective-plane construction~\cite{BRRW2025} gives the lower bound
\[
  C(C_{2k+1},C_{2\ell+1})
  \ge
  \frac{4k^2-1}{4k\ell-1}.
\]
The remaining gap to \eqref{eq:domination-bound} simplifies to
\[
  \frac{2(k-\ell)(2\ell-1)}
  {(4k\ell-1)\bigl(2\ell(2k-2\ell+1)-1\bigr)}.
\]

\section{Formalization blueprint for Lean}
\label{sec:lean}

This section separates the proof into statements whose formal dependencies are
small and explicit.  It is not executable Lean code, but it is intended to
remove mathematical ambiguity before implementation.

\subsection{Recommended finite-dimensional setup}

Use a finite type $V$ with decidable equality and a real matrix
\[
  A:V\to V\to\RR.
\]
For the matrix theorem, assume:
\begin{enumerate}[label=(A\arabic*)]
\item $A$ is symmetric:
  \[
    A_{xy}=A_{yx}\quad\text{for all }x,y.
  \]
\item $A$ is entrywise nonnegative:
  \[
    0\le A_{xy}\quad\text{for all }x,y.
  \]
\end{enumerate}
The exact graph theorem is then obtained by specializing to the adjacency
matrix, whose entries are $0$ or $1$.

A suitable matrix-level statement of the imported theorem is:
\begin{verbatim}
theorem saglam_same_parity
    {V : Type*} [Fintype V] [DecidableEq V]
    (S : Matrix V V Real)
    (hSsymm : S.IsHermitian)
    (hSnonneg : forall i j, 0 <= S i j)
    (u v : V -> Real)
    (hu_nonneg : forall i, 0 <= u i)
    (hv_nonneg : forall i, 0 <= v i)
    (hu_norm : inner u u = 1)
    (hv_norm : inner v v = 1)
    (K R : Nat)
    (hRpos : 0 < R)
    (hRK : R <= K)
    (hparity : K % 2 = R % 2) :
    (inner v ((S ^ K).mulVec u)) ^ R >=
    (inner v ((S ^ R).mulVec u)) ^ K
\end{verbatim}
The exact syntax will depend on the chosen matrix and Euclidean-space APIs.

\subsection{Local lemmas and their proof obligations}

A convenient dependency order is the following.

\begin{enumerate}[label=\textbf{L\arabic*.},leftmargin=2.4em]
\item \emph{Nonnegative powers.}
If $A$ is entrywise nonnegative, then $A^r$ is entrywise nonnegative for every
$r\in\NN$.  Prove by induction using the finite sum defining matrix
multiplication.

\item \emph{Symmetry of powers.}
If $A$ is symmetric, then $A^r$ is symmetric.  This follows from
$(A^r)^\top=(A^\top)^r=A^r$.

\item \emph{Anchor-vector identities.}
For even $b$, define $w_x=A^{b/2}e_x$.  Prove
\[
  \langle w_x,w_x\rangle=(A^b)_{xx}
\]
and
\[
  \langle w_x,A^{r-b}w_x\rangle=(A^r)_{xx}
  \qquad(r\ge b).
\]
The arithmetic obligation is
\[
  \frac b2+(r-b)+\frac b2=r,
\]
which uses that $b$ is even.

\item \emph{Zero anchor case.}
From $\langle w_x,w_x\rangle=0$, conclude $w_x=0$.  Then both later rooted
moments vanish.  This case should be split before normalizing $w_x$ so that no
division by zero occurs.

\item \emph{Positive anchor case.}
Set $u_x=w_x/\sqrt{c_b(x)}$.  Prove coordinatewise nonnegativity and unit norm.
Apply \texttt{saglam\_same\_parity} with
\[
  K=m-b,
  \qquad
  R=n-b.
\]
The parity obligations are
\[
  (m-b)\bmod2=1,
  \qquad
  (n-b)\bmod2=1.
\]
After substitution, clear the positive denominator $c_b(x)^{m-b}$.

\item \emph{Rooted interpolation.}
Package the preceding cases as
\[
  ((A^b)_{xx})^{m-n}((A^m)_{xx})^{n-b}
  \ge((A^n)_{xx})^{m-b}.
\]

\item \emph{Finite two-factor H\"older.}
Formalize Lemma~\ref{lem:integer-holder}.  The cleanest route uses
\texttt{Real.rpow} and the standard finite-sum H\"older inequality with conjugate
exponents
\[
  \frac{m-b}{m-n},
  \qquad
  \frac{m-b}{n-b}.
\]
An alternative is to prove precisely the rational-exponent specialization
needed here and avoid a more general $L^p$ API.

\item \emph{Trace summation.}
Sum the rooted inequality by the H\"older lemma to obtain
\[
  \tr(A^b)^{m-n}\tr(A^m)^{n-b}
  \ge\tr(A^n)^{m-b}.
\]
This matrix statement does not yet require that $A$ be a $0$--$1$ adjacency
matrix.

\item \emph{Trace equals cycle homomorphisms.}
For a simple graph adjacency matrix, expand the finite matrix product and
construct the explicit equivalence between:
\begin{itemize}
\item tuples contributing $1$ to $\tr(A^r)$, and
\item graph homomorphisms $C_r\to T$.
\end{itemize}
Handle $r=2$ separately, proving
\[
  \tr(A^2)=\homn(K_2,T)
\]
from $A_{xy}^2=A_{xy}$.

\item \emph{Density normalization.}
Prove the natural-number identity
\[
  b(m-n)+m(n-b)=n(m-b).
\]
Use it to cancel the common power of $N=|V(T)|$.
\end{enumerate}

\subsection{Suggested main Lean theorem statements}

The matrix statement is the most reusable one:
\begin{verbatim}
theorem trace_even_anchor_odd_interp
    (A : Matrix V V Real)
    (hA_symm : A.IsHermitian)
    (hA_nonneg : forall i j, 0 <= A i j)
    (b n m : Nat)
    (hb_pos : 2 <= b)
    (hbn : b < n)
    (hnm : n < m)
    (hb_even : Even b)
    (hn_odd : Odd n)
    (hm_odd : Odd m) :
    (Matrix.trace (A ^ b)) ^ (m - n) *
    (Matrix.trace (A ^ m)) ^ (n - b) >=
    (Matrix.trace (A ^ n)) ^ (m - b)
\end{verbatim}

The graph statement can then be a corollary:
\begin{verbatim}
theorem cycle_density_even_anchor_odd_interp
    {V : Type*} [Fintype V] [DecidableEq V]
    (T : SimpleGraph V)
    (b n m : Nat)
    (hb_pos : 2 <= b)
    (hbn : b < n)
    (hnm : n < m)
    (hb_even : Even b)
    (hn_odd : Odd n)
    (hm_odd : Odd m) :
    tCycleLike b T ^ (m - n) *
    tCycle m T ^ (n - b) >=
    tCycle n T ^ (m - b)
\end{verbatim}
Here \texttt{tCycleLike 2 T} should be defined as $t(K_2,T)$, while for $b\ge3$ it
is $t(C_b,T)$.

\subsection{Formalization of the domination-exponent consequence}

The exact interpolation theorem is substantially easier to formalize than the
asymptotic domination-exponent corollary.  For the latter, separate the
following components.

\begin{enumerate}[label=\textbf{D\arabic*.},leftmargin=2.4em]
\item The explicit edge-deletion process of
Lemma~\ref{lem:pruning}, including a finite product lower bound.
\item The bijections induced by rotations and reflection of a cycle, used in
\eqref{eq:orientation-symmetry} and
\eqref{eq:incidences-dominate-copies}.
\item The exact gluing bijection
\eqref{eq:gluing-exact}.
\item The spectral proof of
$t(C_{2r},T)\ge t(K_2,T)^{2r}$.
\item The algebraic elimination
\eqref{eq:two-ineq-one}--\eqref{eq:combined-power}.
\item Multiplicativity of homomorphism densities under categorical tensor
products and the real limit in Lemma~\ref{lem:tensor-trick}.
\end{enumerate}

No graphon compactness, cut metric, or graphon approximation theorem is needed
anywhere in this development.

\section{Further structural remarks}

\subsection{Why a vertex Schur complement is not needed here}

The Goodman-style odd-cycle problem and the present domination problem use
regularity in different ways.  In the Goodman-style problem, complementation
is a rank-one perturbation, and the failure of the constant function to be an
eigenfunction creates additive mixed terms.  The constant/mean-zero block
decomposition and its Schur complement are therefore natural tools for
resumming those terms.

Here, after rooting at a vertex, Sa\u{g}lam's theorem gives the pointwise
multiplicative estimate
\[
  \deg_T(x)^{m-n}c_m(x)^{n-2}
  \ge c_n(x)^{m-2}.
\]
The dependence on the nonconstant degree function is already separated into a
single nonnegative factor.  H\"older's inequality performs the exact
regularization
\[
  \sum_x d(x)^\alpha c_m(x)^\beta
  \le
  \left(\sum_xd(x)\right)^\alpha
  \left(\sum_xc_m(x)\right)^\beta.
\]
There is no signed remainder to analyze.  Thus the appropriate general lesson
is not that Schur complementation always removes regularity, but rather:

\begin{quote}
First expose how the regular quantity enters the rooted inequality.  If the
irregularity appears multiplicatively and positively, H\"older may remove it
exactly.  If it appears through additive coupling between the constant and
mean-zero subspaces, a Schur complement is a more natural correction device.
\end{quote}

\subsection{Why the edge-anchored statement is not automatically a graphon inequality}

For a finite simple target, $A_{xy}^2=A_{xy}$, and this is why
\[
  \tr(A^2)=\homn(K_2,T).
\]
For a general graphon $W$,
\[
  \tr(T_W^2)=\int_{[0,1]^2}W(x,y)^2\dd x\dd y,
\]
which is not generally equal to the edge density
$\int W$.  Therefore the direct operator extension of
Theorem~\ref{thm:cycle-interpolation} has the two-edge multigraph density as
its $b=2$ anchor, not $t(K_2,W)$.  The finite simple-target theorem is exactly
the statement needed for the domination exponent.

\subsection{A possible edge-space Schur-complement direction}

The pruning argument controls the distribution of cycle counts rooted at an
edge, not merely the distribution of vertex degrees.  This suggests that any
further improvement beyond Theorem~\ref{thm:domination-bound} should retain
edge-rooted information.  For a graphon $W$ of positive edge density $p$, one
may put the probability measure
\[
  \dd\eta(x,y)=\frac{W(x,y)}p\dd x\dd y
\]
on ordered edges and decompose edge-rooted path kernels into their constant
and mean-zero parts in $L^2(\eta)$.  The gluing identity then contains an exact
covariance term.  A Schur complement on this edge space, rather than on the
original vertex space, is a plausible analogue of the constant/mean-zero
method in the Goodman-style proof.  This observation is only a research
direction and is not used in any theorem above.

\begin{thebibliography}{9}

\bibitem{Saglam2018}
M.~Sa\u{g}lam,
\emph{Near log-convexity of measured heat in (discrete) time and consequences},
arXiv:1808.06717, 2018.

\bibitem{BRwalks}
G.~Blekherman and A.~Raymond,
\emph{A new proof of the Erd\H{o}s--Simonovits conjecture on walks},
Graphs and Combinatorics \textbf{39} (2023), Paper No.~53;
arXiv:2009.10845.

\bibitem{BRRW2025}
G.~Blekherman, A.~Raymond, A.~Razborov, and F.~Wei,
\emph{On domination exponents for pairs of graphs},
arXiv:2506.12151v2, 24 August 2025.

\end{thebibliography}

\end{document}
